\bigskip
\begin{enumerate}
	\item $\forall x \in \R$, $1-\dfrac{1-\e^x}{1+\e^{x}}=\dfrac{1+\e^x}{1+\e^x}-\dfrac{1-\e^x}{1+\e^x}=\dfrac{2\e^x}{1+\e^x}=\dfrac{2\e^x}{\e^x\left(1+\dfrac{1}{\e^x}\right)}=\dfrac{2}{1+\dfrac{1}{\e^x}}=\dfrac{2}{1+\e^{-x}}$

\hfill \fbox{\textbf{Affirmation 1 : Vraie}}
\item $g(x)=\dfrac12 \iff \dfrac{\e^x}{\e^x+1}=\dfrac12 \iff 2\e^x=\e^x+1 \iff \e^x=1 \iff x=0$
\hfill \fbox{\textbf{Affirmation 2 : Vraie}}
	\item La fonction $f$ est continue et dérivable sur $\R$.

	$\forall x \in \R$, $f'(x)=2x\e^{-x}+x^2\times -\e^{-x}= \left(2x-x^2\right)\e^{-x}$.

	L'équation d'une tangente à $\mathcal{C}$ au point d'abscisse $a$ a pour équation $ y=f'(a)(x-a)+f(a)$. Pour que l'axe des abscisses soit tangent à $\mathcal{C}$, il faut que la courbe $\mathcal{C}$ admette une tangente d'équation $y=0$ (équation de l'axe des abscisses), donc il faut que $f(a)=0$ et $f'(a)=0$. Sachant que $\forall a \in \R$, $\e^{-a}\neq 0$,

$f(a)=0 \iff a^2\e^{-a}=0 \iff a^2=0 \iff a=0$. 

$f'(a)=0\iff (2a-a^2)\e^{-a}=0 \iff 2a-a^2=0 \iff a(2-a)=0 \iff a=0 \text{ ou } a=2$.
	
Donc $a=0$. Il n'existe donc qu'un seul point où l'axe des abscisses est tangent à $\mathcal{C}$.

\hfill \fbox{\textbf{Affirmation 3 : Vraie}}
	\item La fonction $h$ est continue et dérivable sur $\R$. Soit $\mathcal{C}_h$ sa courbe représentative.

$\forall x \in \R$, $h'(x)=\e^x\times \left(1-x^2\right)+\e^x\times-2x=\e^x\left(1-x^2-2x\right)= \left(-x^2-2x+1\right)~\e^x$.

La fonction $h'$ est continue et dérivable sur $\R$.

$\forall x \in \R$, $h''(x)=(-2x-2)\e^x+\e^x\times\left(-x^2-2x+1\right) =\e^x\left(-2x-2-x^2-2x+1\right)= \left(-x^2-4x-1\right)\e^x$

Si $\mathcal{C}_h$ admet des points d'inflexions, alors $h''(x)$ peut s'annuler et changer de signe. Or $\e^x > 0$ pour tout réel $x$, donc étudions le signe de $- x^2- 4x- 1$ sur $\R$.

$-x^2-4x-1$ s'annule pour $x_1=-2-\sqrt{3}$ et pour $x_2=-2+\sqrt{3}$ car $\Delta = 12$.

Le trinôme du second degré s'annule donc deux fois et change donc aussi de signe.

\hfill \fbox{\textbf{Affirmation 4 : Fausse}}
	\item $\forall x \in \R$, $f(x)=\dfrac{\e^x}{\e^x+x}=\dfrac{\e^x}{\e^x\left(1+\dfrac{x}{\e^x}\right)}=\dfrac{1}{1+\dfrac{x}{\e^x}}$.

D'après les croissances comparées, $\lim \limits_{n \rightarrow +\infty} \dfrac{\e^x}{x}=+\infty$ donc $\lim \limits_{n \rightarrow +\infty} \dfrac{x}{\e^x}=0$ donc $\lim \limits_{n \rightarrow +\infty} f(x)=1$

\hfill \fbox{\textbf{Affirmation 5 : Fausse}}
		\item $1+\e^{2x}\geqslant 2\e^x \iff \e^{2x}-2\e^x+1\geqslant0 \iff \left(\e^x-1\right)^2\geqslant0$.

Cette dernière inégalité est vraie pour tout réel $x$.
\hfill \fbox{\textbf{Affirmation 6 : Vraie}}
\end{enumerate}
